\section{Receta de extremo a extremo: la lógica de composición de OLMo a Tulu}

\subsection{Principio de composición: primero estabilizar la base, luego hacer la alineación y al final optimizar el presupuesto}
La ruta práctica real que propone esta clase no es «elegir primero el algoritmo», sino «elegir primero el orden de composición». Un orden ejecutable es:
\begin{enumerate}
    \item En la etapa de preentrenamiento, primero se establece una base estable;
    \item SFT ajusta el protocolo de la tarea y el comportamiento de salida a un rango controlable;
    \item DPO/PPO realiza el ajuste fino de preferencias;
    \item RLVR impulsa el razonamiento en subdominios verificables;
    \item el escalado en tiempo de prueba atiende las solicitudes de cola larga de alto valor.
\end{enumerate}

\begin{knowledgebox}{Por qué esta ruta es favorable para la comunidad abierta}
Cada capa puede experimentarse, evaluarse y liberarse como código abierto de forma independiente, lo que reduce la barrera de entrada para la colaboración en investigación y facilita localizar el problema cuando falla la reproducción de un experimento.
\end{knowledgebox}

\subsection{Sistema de evaluación: no basta con mirar un solo benchmark}
\begin{importantbox}{La perspectiva de evaluación de la clase}
La mejora de Reasoning debe verificarse conjuntamente en varias dimensiones:
\begin{itemize}
    \item tareas verificables como matemáticas y código;
    \item preguntas y respuestas de conocimiento y exactitud factual;
    \item estabilidad en contextos largos;
    \item restricciones de multilingüismo y seguridad;
    \item métricas de costo y latencia.
\end{itemize}
\end{importantbox}

\begin{warningbox}{El éxito en las métricas no equivale al éxito del sistema}
Si la evaluación no cubre las restricciones de despliegue (latencia, presupuesto, robustez), el modelo puede resultar inutilizable en un entorno de producto real aunque mejore en el benchmark.
\end{warningbox}

\subsection{Resumen de la sección}
El núcleo de la receta de extremo a extremo es el orden y la composición, no la optimización de un punto aislado. Los sistemas exitosos suelen ser el resultado de múltiples etapas «subóptimas pero complementarias».


